\section{Considerações finais}
\label{sec:consideracoes}

A \projeto{} está na fase de validação e dispõe de um núcleo de produto funcional, evidências qualitativas favoráveis e um modelo econômico explícito. O diagnóstico mostra que o principal obstáculo ao crescimento não é a falta de usuários em potencial, dado o grande número de lares com cães e gatos \cite{correio2020}, mas a baixa contribuição por usuário e a dependência de convênios.

Por essa razão, as duas estratégias propostas atuam de forma complementar: a densidade local por meio de parcerias cria a base de tutores e clínicas, e o aprofundamento de receita eleva a contribuição de cada usuário e reduz o custo do fundo de créditos. As ações de escalabilidade, centradas em automação, infraestrutura segura, pacote replicável por praça e receita recorrente, buscam desvincular o crescimento do esforço manual da equipe. A simulação indica que as alavancas combinadas elevam a contribuição de R\$~0,55 para R\$~1,42 por usuário e reduzem de cerca de 10.155 para cerca de 3.918 os usuários necessários ao equilíbrio do cenário com equipe remunerada.

O plano é, assim, uma hipótese estruturada e verificável, acompanhada de metas, gatilhos e condições de revisão. O piloto em Belo Horizonte deve produzir os dados que faltam (CAC, retenção, conversão e disposição a pagar), e o grupo deve atualizar este documento e a simulação com esses dados antes de qualquer decisão de expansão.
